\section{Apéndice}

\subsection{Entorno de ejecución}

La Tabla~\ref{tab:entorno} detalla las herramientas y versiones utilizadas en la implementación experimental.

\begin{table}[ht]
\centering
\caption{Herramientas de software utilizadas.}
\label{tab:entorno}
\begin{tabular}{lll}
\hline
\textbf{Herramienta} & \textbf{Versión} & \textbf{Propósito} \\
\hline
Python \cite{python3} & 3.12+ & Lenguaje de programación \\
Pyomo \cite{pyomo2021} & 6.7+ & Modelado algebraico de optimización \\
HiGHS \cite{highs2024} & 1.7+ & Solver LP/MIP \\
NetworkX \cite{networkx2008} & 3.2+ & Manipulación de grafos \\
NumPy \cite{numpy2020} & 1.26+ & Cómputo numérico \\
Matplotlib \cite{matplotlib2007} & 3.8+ & Visualización \\
Jupyter \cite{jupyter2016} & 7.0+ & Notebook interactivo \\
\hline
\end{tabular}
\end{table}

\subsection{Fragmento de código: modelo ILP en Pyomo}

A continuación se presenta el código principal de la formulación del MCP en Pyomo:

\begin{verbatim}
import pyomo.environ as pyo

def build_mcp_model(G):
    model = pyo.ConcreteModel("MaximumClique")
    model.V = pyo.Set(initialize=list(G.nodes()))
    model.x = pyo.Var(model.V, domain=pyo.Binary)
    model.obj = pyo.Objective(
        expr=sum(model.x[i] for i in model.V),
        sense=pyo.maximize
    )
    model.nonadj = pyo.ConstraintList()
    for i in G.nodes():
        for j in G.nodes():
            if i < j and not G.has_edge(i, j):
                model.nonadj.add(
                    model.x[i] + model.x[j] <= 1
                )
    return model

model = build_mcp_model(G)
solver = pyo.SolverFactory('appsi_highs')
result = solver.solve(model)
\end{verbatim}

\subsection{Instancia keller4: formato DIMACS}

La instancia se distribuye en formato ASCII estándar DIMACS con la siguiente estructura:

\begin{verbatim}
c keller4.clq
c Keller graph on 4 dimensions
p col 171 9435
e 1 2
e 1 3
...
\end{verbatim}

Donde \texttt{p col 171 9435} define un grafo de 171 vértices y 9\,435 aristas, y cada línea \texttt{e u v} define una arista entre los vértices $u$ y $v$. El archivo completo está disponible en el repositorio del proyecto.